% ====================================================================
% Capítulo 1: Introducción
% ====================================================================
\chapter{Introducción y panorama histórico}
\label{ch:introduccion}

\epigraph{\emph{``Si un fluido caliente se expone al frío, se enfría más rápido que un fluido frío.''}}
{--- Aristóteles, \emph{Meteorología} I.12 (350 a.C.)}

\section{El fenómeno en su definición operacional}

Se denomina \textbf{efecto Mpemba} al fenómeno termodinámico mediante el cual un 
sistema preparado a una temperatura inicial alta $\Th$ alcanza el equilibrio térmico 
con un baño a temperatura $\Tb < \Th$ en un tiempo menor que un sistema idéntico 
preparado a una temperatura inicial intermedia $\Tw$, con $\Tb < \Tw < \Th$. 
Formalmente, sean $p(t;\Tinit)$ la distribución de probabilidad sobre el espacio de 
estados al tiempo $t$ dado que el sistema fue preparado en equilibrio térmico a 
$\Tinit$ y luego acoplado al baño $\Tb$, y sea $D[\cdot,\cdot]$ una función de 
distancia al equilibrio (verificando $D[p,\pieq] \ge 0$ con igualdad si y solo si 
$p = \pieq$). El efecto Mpemba ocurre cuando
\begin{equation}
    \exists\, t^{\ast} > 0 \;\text{tal que}\; 
    D\!\left[p(t^{\ast};\Th),\pieq(\Tb)\right] 
    < D\!\left[p(t^{\ast};\Tw),\pieq(\Tb)\right]
    \label{eq:mpemba_def}
\end{equation}
y la desigualdad persiste para todo $t > t^{\ast}$~\cite{lu2017}.

El término $t^{\ast}$ recibe el nombre de \emph{tiempo de cruzamiento} o 
\emph{crossover time}. La condición \cref{eq:mpemba_def} parece contraintuitiva: si la 
relajación térmica fuera estrictamente cuasi-estática, el sistema más caliente 
debería pasar por todos los estados intermedios y nunca podría ``adelantar'' a un 
sistema que partió más cerca del equilibrio. La paradoja se resuelve cuando se 
reconoce que la dinámica no-equilibrio recorre trayectorias en el espacio de 
probabilidades cuyo arco geodésico depende fuertemente de las amplitudes iniciales de 
los modos de decaimiento espectral del operador de evolución temporal.

\section{Notas históricas: 2350 años de observación}

La afirmación de que el agua caliente puede congelarse más rápido que el agua fría 
aparece registrada en la \emph{Meteorología} de Aristóteles, quien describe la 
práctica de los habitantes del Ponto de exponer el agua al sol antes de hacer agujeros 
en el hielo, con la creencia de que así se solidificaba más rápido. Francis Bacon 
en el \emph{Novum Organum} (1620) y René Descartes en \emph{Les Météores} (1637) 
mencionaron observaciones análogas, sin proponer mecanismo físico alguno~\cite{burridge2016, hallstadius2020}.

Joseph Black, en 1775, fue el primer científico moderno en investigar el fenómeno: 
comparó el comportamiento de muestras de agua hervida frente a muestras sin hervir y 
observó que la eliminación de gases disueltos y la agitación mecánica modificaban 
drásticamente la dinámica de cristalización. Sin embargo, durante casi dos siglos 
el efecto careció de tratamiento sistemático, y fue desestimado como artefacto 
experimental.

\section{El redescubrimiento de Mpemba (1963--1969)}

La denominación moderna del fenómeno proviene de un estudiante de secundaria 
tanzano, Erasto B. Mpemba, quien en 1963 observó durante una clase de cocina que 
una mezcla de helado preparada con leche caliente se solidificaba en el congelador 
antes que una preparada con leche fría. Tras años de escepticismo por parte de sus 
profesores, Mpemba colaboró con Denis G. Osborne, quien publicó en 1969 el primer 
estudio sistemático en \emph{Physics Education}~\cite{mpemba1969}. Sus datos más 
notables: agua inicialmente a $\SI{80}{\celsius}$ comenzó a congelarse en menos 
de un tercio del tiempo que el agua inicialmente a $\SI{20}{\celsius}$ en condiciones 
nominalmente idénticas.

\section{El problema de la reproducibilidad y la definición rigurosa}

A pesar de múltiples reportes confirmando el efecto en agua~\cite{burridge2016}, las 
discrepancias entre experimentos plantearon una duda fundamental: \emph{¿es el 
efecto Mpemba un fenómeno físico genuino o un artefacto de las condiciones no 
controladas de los experimentos macroscópicos?}

Burridge y Linden~\cite{burridge2016} realizaron en 2016 un meta-análisis crítico de 
los experimentos clásicos y concluyeron que, bajo la definición estricta de 
``condiciones idénticas'' (mismo recipiente, masa, configuración térmica, ausencia 
de sitios de nucleación diferenciales), \emph{no existía evidencia estadísticamente 
significativa} del efecto Mpemba en agua. Su conclusión: el efecto reportado 
históricamente derivaba en gran medida de la varianza estocástica intrínseca a la 
nucleación cristalina heterogénea.

Esta conclusión espoleó la búsqueda de protocolos experimentales rigurosos. En 
particular, Hallstadius y Burridge~\cite{hallstadius2020} demostraron en 2020 que 
\emph{sí} es posible observar el efecto Mpemba de manera repetible cuando se 
controla \emph{sistemáticamente} la asimetría de sitios de nucleación entre las 
muestras: específicamente, lijando con papel de esmeril el interior del recipiente 
que contiene la muestra caliente.

\section{La reformulación moderna del fenómeno}

A partir de 2017, la formalización matemática del efecto Mpemba en sistemas de 
relajación markoviana ha permitido superar el debate clásico centrado exclusivamente 
en el agua. Las contribuciones clave son:

\begin{itemize}
\item \textbf{Lu \& Raz (2017)}~\cite{lu2017}: formalizan el efecto en sistemas 
markovianos genéricos y derivan una condición suficiente para su aparición en 
términos del espectro de la matriz de tasas. Predicen además el \emph{efecto Mpemba 
inverso} (calentamiento anómalo).
\item \textbf{Klich, Raz, Hirschberg \& Vucelja (2019)}~\cite{klichraz2019}: 
introducen el concepto de \emph{efecto Mpemba fuerte} ---la cancelación geométrica 
del modo de decaimiento más lento--- y un índice topológico que cuenta el número 
de tales preparaciones óptimas.
\item \textbf{Lasanta \emph{et al.} (2017)}~\cite{lasanta2017}: prueban la 
aparición del efecto en gases granulares mediante una formulación de momentos 
truncada en el coeficiente de Sonine.
\item \textbf{Kumar \& Bechhoefer (2020)}~\cite{kumar2020}: realizan la primera 
verificación experimental \emph{controlada} del efecto Mpemba fuerte, en un 
sistema coloidal sometido a una pinza óptica con potencial efectivo programable.
\item \textbf{Baity-Jesi \emph{et al.} (2019)}~\cite{baityjesi2019}: muestran 
mediante simulaciones a gran escala (Janus II) que el efecto Mpemba es una 
manifestación robusta de la memoria estructural en vidrios de espín.
\item \textbf{Carollo, Lasanta \& Lesanovsky (2021)}~\cite{carollo2021}: extienden 
el formalismo al régimen cuántico abierto markoviano mediante la ecuación 
maestra de Gorini--Kossakowski--Sudarshan--Lindblad.
\item \textbf{Vu \& Hayakawa (2025)}~\cite{vuhayakawa2025}: unifican 
matemáticamente el fenómeno mediante el preorden de termomayorización, 
estableciendo una definición independiente de la métrica de distancia 
seleccionada.
\item \textbf{Chattopadhyay, Santos \& Misra (2026)}~\cite{chattopadhyay2026} y 
\textbf{Li \& Yang (2026)}~\cite{liyang2026}: convierten el efecto Mpemba en un 
recurso metrológico para termometría cuántica de no-equilibrio.
\end{itemize}

\section{Estructura de este documento}

Los capítulos siguientes desarrollan cada uno de estos marcos con detalle teórico 
y verificación numérica. La organización es:

\begin{description}
\item[Capítulo \ref{ch:debate_clasico}.] El debate experimental clásico en agua, 
los criterios de Burridge--Linden y el protocolo de nucleación controlada.
\item[Capítulo \ref{ch:mecanismos_agua}.] Mecanismos moleculares: enlaces de 
hidrógeno fuertes, supersolidez de piel, hexámeros de agua y el paradigma 
polimérico de Naserifar--Goddard.
\item[Capítulo \ref{ch:markoviano}.] Formulación markoviana de Lu \& Raz, 
condiciones espectrales y su versión inversa.
\item[Capítulo \ref{ch:mpemba_fuerte}.] El efecto Mpemba fuerte, el modelo de 
energía aleatoria (REM) y el índice topológico de Klich--Raz.
\item[Capítulo \ref{ch:granular}.] Fluidos granulares: ecuaciones de momentos, 
DSMC y mezclas binarias.
\item[Capítulo \ref{ch:coloides}.] Verificación experimental en coloides 
brownianos: efecto directo (Kumar 2020) e inverso (Kumar--Chétrite 2022).
\item[Capítulo \ref{ch:spin_glass}.] El efecto Mpemba en vidrios de espín como 
memoria persistente: la simulación Janus II y la correlación entre longitud de 
coherencia y energía interna.
\item[Capítulo \ref{ch:cuantico}.] El régimen cuántico: ecuación de Lindblad, 
modos de Liouville, puntos excepcionales y el efecto Mpemba metrológico.
\item[Capítulo \ref{ch:termomajorizacion}.] La unificación de Vu--Hayakawa 
mediante el preorden de termomayorización.
\item[Capítulo \ref{ch:metodos_numericos}.] Métodos numéricos en detalle: 
integradores Runge--Kutta de cuarto orden, DSMC, dinámica de Langevin estocástica, 
descomposición espectral por rotaciones de Jacobi y dinámica molecular 
gruesa-grano.
\item[Capítulo \ref{ch:implementacion}.] La suite \texttt{Mpemba-X v2}: 
arquitectura del código, organización de los módulos, paralelización OpenMP y 
MPI, y formato de los archivos de configuración.
\item[Capítulo \ref{ch:conclusiones}.] Síntesis y direcciones futuras.
\end{description}

Las pruebas matemáticas técnicas, los \emph{listings} completos del código y las 
demostraciones del operador de termomayorización se encuentran en los apéndices.
